\section{The matrix construction target}
\label{sec:convex-bodies}

The geometric part of the argument is an elementary conversion of a real
matrix into a pair of convex bodies. It specifies exactly what the two
combinatorial constructions must achieve: many separated rows and a small
uniform approximation list for the absolutely convex hull of the columns.
The small cube added to the row hull ensures nonempty interior without
losing the polar covering estimate.

\begin{lemma}[Matrix-to-body conversion]
\label{lem:matrix-to-bodies}
Let $X$ and $Y$ be finite nonempty sets, and let
$g_y\colon X\to\R$, $y\in Y$, be real functions.  Write
\[
 R_x=(g_y(x))_{y\in Y}\in\R^Y,
 \qquad
 f_\lambda=\sum_{y\in Y}\lambda_y g_y,
 \qquad
 \mathcal F=\{f_\lambda:\|\lambda\|_1\le1\}.
\]
Suppose that
\[
 \|R_x-R_{x'}\|_\infty\ge1
 \quad\text{for all distinct }x,x'\in X,
\]
and that, for some $\varepsilon>0$, there is a list of $M$ functions on
$X$ such that every member of $\mathcal F$ is within
$\varepsilon$ of a function on the list in the uniform norm.
The functions on the list need not belong to $\mathcal F$.
For $t>0$, set
\[
 K=3\absconv\{R_x:x\in X\}+tB_\infty^Y,
 \qquad L=B_\infty^Y.
\]
Then $K,L$ are origin-symmetric compact convex bodies with nonempty
interior in $\R^Y$, and
\begin{equation}
\label{eq:matrix-body-bounds}
 N(K,L)\ge |X|,
 \qquad
 N\bigl(L^\circ,(6\varepsilon+2t)K^\circ\bigr)\le M.
\end{equation}
Both covering numbers use arbitrary ambient translation centers.
\end{lemma}

\begin{proof}
The hull $\absconv\{R_x:x\in X\}$ is the convex hull of the
finitely many vectors $\pm R_x$.  It is compact, convex, and symmetric,
and contains zero.  Thus $K$ has the same three properties and contains
$tB_\infty^Y$, which gives it nonempty interior in the full ambient
space.  The corresponding properties of $L$ are immediate.

Each point $3R_x$ belongs to $K$, and two distinct such points have
uniform distance at least $3$.  Every translate of $L$ has uniform
diameter $2$, whatever its center.  Hence it contains at most one of
these $|X|$ points.  This proves the first inequality in
\eqref{eq:matrix-body-bounds}.

For the polar estimate, the support function
$h_K(\delta)=\sup_{z\in K}\langle z,\delta\rangle$ is exactly
\begin{equation}
\label{eq:row-support-function}
 h_K(\delta)
 =3\max_{x\in X}|\langle R_x,\delta\rangle|
     +t\sum_{y\in Y}|\delta_y|
 =3\|f_\delta\|_{\infty,X}+t\|\delta\|_1.
\end{equation}
Indeed, a linear functional attains its maximum over the finite convex
hull at one of the vectors $\pm R_x$.  Its maximum over
$B_\infty^Y$ is $\sum_y|\delta_y|$, attained by choosing the sign of
each coordinate.  The two choices in the Minkowski sum are independent,
so their maxima add.  The same cube calculation gives
\[
 L^\circ=B_1^Y=\{\lambda\in\R^Y:\|\lambda\|_1\le1\}.
\]

Order the approximating list.  Assign each $\lambda\in B_1^Y$ to the
first list function within $\varepsilon$ of $f_\lambda$, thereby
partitioning $B_1^Y$ into at most $M$ nonempty clusters.  From each
nonempty cluster choose one actual vector $\lambda^{(k)}\in B_1^Y$.
For any $\lambda$ in that cluster, the triangle inequality gives
\[
 \|f_\lambda-f_{\lambda^{(k)}}\|_{\infty,X}
 \le2\varepsilon,
 \qquad
 \|\lambda-\lambda^{(k)}\|_1\le2.
\]
Consequently \eqref{eq:row-support-function} implies
\[
 h_K(\lambda-\lambda^{(k)})\le6\varepsilon+2t.
\]
For every $c>0$, the definition of the polar and positive homogeneity
give the equivalence
\[
 \delta\in cK^\circ
 \quad\Longleftrightarrow\quad h_K(\delta)\le c.
\]
Each cluster is therefore contained in
$\lambda^{(k)}+(6\varepsilon+2t)K^\circ$.  These at most $M$
translates cover $B_1^Y=L^\circ$, proving the second inequality.
Selecting representatives inside the coefficient ball is what controls
the contribution of the added cube, even when the original list
functions lie outside $\mathcal F$.
\end{proof}

For proposed constants $a,b\ge1$, the lemma reduces the problem to finding
a matrix and positive numbers $\varepsilon,t$ for which
\[
 \log|X|>b\log M,
 \qquad 6\varepsilon+2t\le a^{-1}.
\]
Indeed, increasing the polar covering body preserves the cover supplied
by the lemma. The next two sections construct the matrix. The compression
argument controls $M$ through the number of labels in a partition; the
finite-field construction makes that number small relative to the row
count on the logarithmic scale.
